\section{A stable Hecke component inside the exact tame character}
\label{sec:cubic-hsieh-tame-stable}

Retain the imaginary quadratic setting, original prime-to-$p$
branch $\beta$, all published arithmetic inputs, and the exact
integral constructions of AC-G67-measure, AC-G67-local,
AC-G67-boundary and AC-G67-type.  Thus $p>3$ splits in $K$,
$p\nmid h_K$, the branch is unramified above $p$, and
\[
 \beta'=\beta\vartheta,
 \qquad \Lambda'=\mathcal O'[[\Gamma_K]]
             \simeq\mathcal O'[[X,Y]],
 \qquad \lambda=(\vartheta_q^{-1},1,\vartheta_q),
                                                        \tag{G70.1}
\]
where $\vartheta$ is the rational tame character of order
$p^r$ at the first added split prime $q$.
The coefficient ring contains its values.  Assume the
residual condition \textup{(Dist)} of Hsieh (8.2).
The following result concerns this enlarged ordinary Hsieh
family.  It does not identify that family with the actual
symmetric-weight deformation problem.

\begin{proposition}[AC-G70: a generically irreducible component with the prescribed tame character]
\label{AC-G70}
One can add a second split prime $q'$, outside the preceding
fixed finite set, and carry out the G67 construction at the
same common level for $\beta$ and $\beta'$, so that the
following assertions hold.

The completed algebra generated by the away-level Hecke
operators on the integral ordinary cuspidal
$\lambda$-character submodule is finite, reduced and
torsion-free over $\Lambda'$.  Its Eisenstein local factor
$R_\lambda$ is nonzero.  If $I_\lambda$ is its actual
Eisenstein ideal, then
\[
 R_\lambda/I_\lambda\simeq\Lambda'/\operatorname{Eis}_\lambda,
 \qquad
 \operatorname{Eis}_\lambda\subset(L_D^S L_K^S).
                                                        \tag{G70.2}
\]
Its residual pseudocharacter is precisely
\[
           \bar\beta^{-c}\bar\kappa_p\;\oplus\;1
                        \;\oplus\;\bar\beta\bar\kappa_p^{-1}.
                                                        \tag{G70.3}
\]
There is a height-one prime $P\subset\Lambda'$ with
$\operatorname{ord}_P L_K^S>0$ and
$\operatorname{ord}_P L_D^S=0$, and a prime
$\mathfrak q\subset R_\lambda$ over $P$ containing
$I_\lambda$.  Every minimal component through
$\mathfrak q$ has generically absolutely irreducible
three-dimensional pseudocharacter and contains classical
ordinary cuspidal points with absolutely irreducible attached
Galois representation.  Their automorphic representation at
$q$ contains the exact character (G70.1), hence has the
principal-series local parameter of AC-G67-type.

No local type is asserted at the second prime $q'$.
At $q$ the passage from that automorphic parameter to the
attached Galois inertia still requires the precise local--global
compatibility statement retained in AC-G67-type.
The repeated-inertia issue at $269$, the symmetric ordered
Hodge type, and the actual Hecke-support containment of
AC-G61 remain separate requirements.
\end{proposition}
\begin{proof}
\proofstep{Choose the second prime before fixing the common integral family.}
Let $L$ be the normal closure over $\mathbb Q$ of the field
cut out by $\beta'$ together with $K(\mu_p)$.
Choose $q'$ completely split in $L$, outside the original
bad places and $q$.  The elementary prime-divisor argument
in AC-K29-Hsieh-aux applies to any finite normal extension;
it does not require $p\nmid[L:K]$.
Thus $\beta'(\operatorname{Frob}_{w'})=1$ for both
$w'\mid q'$.

First apply the original prime-to-$p$ nonvanishing input to
$\beta$, with both $q$ and $q'$ prescribed in the split
level product.  The proof of Hsieh Proposition~7.13 excludes
only finitely many auxiliary-character values at each
additional fixed prime; its repaired simultaneous
nonvanishing input supplies an infinite set.  There are
only finitely many added primes, so the same proof supplies
a primitive Fourier--Jacobi functional at this common level.
The fixed-level measure equality of AC-G67-measure then
transfers primitivity to $\beta'$.  AC-G67-local and
AC-G67-boundary impose the exact $q$-character and produce
the primitive cuspidal congruence of AC-G67-type.
No second application of a prime-to-$p$ branch-field
nonvanishing theorem to $\beta'$ has been made.
No assertion about the type at $q'$ is part of this step.

\proofstep{The integral character algebra is finite and is already complete.}
At the pro-unipotent $q$-Iwahori level used in G67, let $M$
be the full ordinary module and let $M^0$ be its cuspidal
submodule.  Hsieh Theorem~4.21, including its full-module
case, and the finite coefficient extension give a finite
free ambient $\Lambda'$-module $M$.
Put
\[
 M^0_\lambda=\{m\in M^0:a m=\lambda(a)m\text{ for every }a\in A\},
 \qquad A=(\mathbb F_q^\times)^3.
                                                        \tag{G70.4}
\]
It is a finite torsion-free module: it is a submodule of
$M$ over the Noetherian domain $\Lambda'$.
We do not claim that it is integrally free or an integral
direct summand.  After inverting $p$,
\[
 M^0_\lambda[1/p]=e_\lambda M^0[1/p],
 \qquad e_\lambda=|A|^{-1}\sum_{a\in A}\lambda(a)^{-1}a.
                                                        \tag{G70.5}
\]
All other factors of $|A|$ are units already.
This identity is used only after inverting $p$; the integral
congruence was obtained by the trace-one lift in G67.

Let $h_\lambda$ be the $\Lambda'$-algebra generated by the
away-$S\cup S_p$ Hecke operators in
$\operatorname{End}_{\Lambda'}(M^0_\lambda)$.
This endomorphism module is finite: a finite presentation of
$M^0_\lambda$ embeds it in a finite direct sum of copies of
$M^0_\lambda$.  Every $\Lambda'$-submodule of a finite
module is finite.  Therefore $h_\lambda$ is finite and
torsion-free over $\Lambda'$.
It is complete and closed in the endomorphism module for
the maximal-ideal topology: both it and the quotient are
finite modules over a complete Noetherian local ring, and
the quotient is separated.  Consequently taking the
completed Hecke algebra adds no elements and changes no
relations.  Equation (G70.5) likewise makes
$h_\lambda[1/p]$ a finite $\Lambda'[1/p]$-algebra, not an
uncontrolled completion of a rational character projector.

\proofstep{Classical control, reducedness, and density hold on this exact character.}
The finite group $A$ acts at a prime different from $p$;
its action commutes with ordinary control and with every
away-level Hecke operator.  After inverting $p$, applying
$e_\lambda$ to the control isomorphism therefore gives
exactly the $\lambda$-subspace of classical cusp forms.
The usual semisimplicity of the commuting classical
unramified Hecke operators, used in Hsieh
Proposition~7.17, applies to this invariant subspace as well.

Here is also an explicit check that no component is lost
in passing from the full module to (G70.4).
Over $Q=\operatorname{Frac}(\Lambda')$, choose bases for
$M^0_\lambda\otimes Q$ and $h_\lambda\otimes Q$.
Clear their finitely many denominators.  Then, on some
nonempty principal open of $\operatorname{Spec}\Lambda'$, both
modules are free.  The faithful inclusion
$h_\lambda\hookrightarrow\operatorname{End}(M^0_\lambda)$
has a nonzero maximal minor over $Q$; invert that minor
too.  Every remaining specialization is consequently a
faithful matrix representation of the specialized algebra.
At a classical weight in this open, its image is the
actual semisimple Hecke algebra acting on the classical
$\lambda$-space.  All its points are classical eigenforms.

A nilpotent in $h_\lambda$ acts by zero at every such
classical specialization.  Classical weights in the
separable cone are Zariski dense in weight space, so all
matrix coefficients vanish on the chosen open.  The
endomorphism is zero over $Q$ and hence zero integrally
by torsion-freeness.  Thus $h_\lambda$ is reduced.
Each minimal prime contracts to zero in $\Lambda'$:
a nonzero element of the base lying in a minimal prime
of a reduced Noetherian ring is a zero divisor, contrary
to base torsion-freeness.  Hence every component dominates
weight space.  Classical points are dense on each one:
for a proper closed subset of a component, its finite
image in weight space is proper, by the dimension theorem
for finite integral extensions; the dense classical
weights avoiding that image and the preceding principal
open lift to classical points outside the subset.
This argument also allows arbitrarily regular classical
weights in the same source cone.

\proofstep{The integral congruence selects the same residual constituent.}
The primitive cuspidal congruence of G67 gives
\[
 h_\lambda/I_\lambda\twoheadrightarrow
                       \Lambda'/(L_D^S L_K^S).
                                                        \tag{G70.6}
\]
Modulo the ideal generated by the differences between the
away-level Hecke operators and their prescribed Eisenstein
eigenvalues, the algebra is generated by scalars.
Thus its quotient is exactly
$\Lambda'/\operatorname{Eis}_\lambda$, with
$\operatorname{Eis}_\lambda\subset(L_D^SL_K^S)$.
A finite algebra over a complete local ring is a product
of its complete local factors.  The unique factor through
which this nonzero local scalar quotient passes is
$R_\lambda$; all other factors have unit Eisenstein ideal.
This proves (G70.2) and preserves all the properties just
proved.

At every classical point the same Galois representations
as in Hsieh \S8.3 supply a pseudocharacter.  One may obtain
it by the source's dense classical patching, or by restricting
the full-level pseudocharacter to this Hecke image; equality
at dense classical points verifies the quotient relations.
It is continuous, conjugate self-dual, and modulo
$I_\lambda$ equals
\[
 \delta_1+\delta_2+\delta_3,
 \qquad
 \delta_1=\Psi'{}^{-c}\kappa_p,\quad\delta_2=1,
 \quad\delta_3=\Psi'\kappa_p^{-1},
 \quad\Psi'=\beta'\varepsilon_{\Lambda'}.
                                                        \tag{G70.7}
\]
Indeed these equalities hold on the away-level Frobenius
set and then on the Galois group by continuity and
Chebotarev.  Since $\vartheta\equiv1$ modulo the coefficient
maximal ideal, reduction is exactly (G70.3).
Enlarging $S$ by $q'$ deletes only finitely many Frobenius
conditions and changes neither that conclusion nor
\textup{(Dist)}.

\proofstep{An Euler divisor away from the cyclotomic and ordinary coincidences.}
For $w'\mid q'$, its imprimitive factor is
\[
 f_{w'}=1-[\operatorname{Frob}_{w'}]
       =1-(1+X)^a(1+Y)^b,\qquad b\ne0,
                                                        \tag{G70.8}
\]
where $X,Y$ are cyclotomic and anticyclotomic coordinates.
The full proof of AC-G64-prime applies over $\mathcal O'$:
if $w'{}^{h_K}=(\pi)$, vanishing of the anticyclotomic
logarithm would make $\pi/\bar\pi$ a root of unity,
contradicting its nontrivial ideal.  No prime factor of
$f_{w'}$ divides a nonzero cyclotomic-only element.
Choose such a prime factor $P$.  The integral imprimitive
factorization gives $P\mid L_K^S$, while $P\nmid L_D^S$
and $p\notin P$.

The local ordinary separation used in the following step
can also be checked directly for this divisor.
The pro-$p$ ray-class sequence, using $p\nmid h_K$ and the
finite prime-to-$p$ global unit group, identifies
$\Gamma_K$ with the two local principal-unit directions
at $\mathfrak p,\bar{\mathfrak p}\mid p$.
Use their coordinates $U,V$.
The two logarithmic coordinates of the split Frobenius
are nonzero: up to fixed nonzero factors they are
$\log_p\iota_p(\pi)$ and
$\log_p\iota_p(\bar\pi)$.  Each argument is a $p$-adic
unit, and either logarithm could vanish only if the
corresponding algebraic number were a root of unity,
contrary to its nontrivial ideal.
Thus (G70.8), expressed in $U,V$, has nonzero exponents
in both directions.  It has no prime factor depending
only on $U$ or only on $V$: after reduction by such a
prime, the other power series is still nonconstant;
in residue characteristic $p$ its first nonconstant
coefficient occurs in a $p$-power degree and is nonzero.
Since $\beta'$ is unramified at $p$, the equations
$\delta_3=\delta_2$ and $\delta_1=\delta_2$ on a local
inertia group are respectively single-direction equations
of this kind, with a fixed cyclotomic constant.
The equation $\delta_1=\delta_3$ is cyclotomic-only.
Allowing a finite index in the chosen principal units
only replaces these equations by their $p$-power versions,
whose prime factors have the same respective property.
None can hold modulo $P$.  This verifies the exact
inertia distinction required in Hsieh Lemma~8.7.

By (G70.2) there is a prime $\mathfrak q$ over $P$
containing $I_\lambda$.  Fix a minimal prime
$\mathfrak a\subset\mathfrak q$.

\proofstep{The faithful-trace argument retains this character submodule.}
We specify precisely the published inputs being reused.
Hsieh \S8.3, (8.2)--(8.9), constructs the residually
multiplicity-free conjugate-dual generalized matrix
algebra from the faithful trace quotient of
Bella\"iche--Chenevier Theorem~1.4.4(ii).
The construction uses the pseudocharacter, the finite
reduced Hecke algebra, and its dense classical points;
all have just been verified for $R_\lambda$.
The ordinary local calculation in \S8.4 uses the same
classical $p$-local type; the tame projector at $q$ changes
none of it.  Corollary~8.9 supplies the injection of the
first off-diagonal quotient into the indicated Selmer
group.  Lemma~8.10 uses only quadratic restriction with
$p>2$.  The Wiles inequality in Theorem~8.11 is stated
for the finite branch character in question and contains
no hypothesis $p\nmid\#\Delta$.
The separate ray-class identification (8.1) in \S8.2.3
\emph{does} have that hypothesis and is not used here.
The hypotheses (1)--(2) of Theorem~7.18 retained at the
start of \S8 hold: the field and coefficient prime are
unchanged and $\beta'$ remains unramified above $p$.
Primitivity, rather than a new branch-field nonvanishing
assertion, was proved in the first step.

For completeness, the remaining algebra is as follows.
In the faithful trace quotient write the finite fractional
entry modules as $R_{ij}$.  Conjugate duality gives
$R_{12}=R_{23}$, and put
\[
 I_1=R_{13},\qquad J_1=R_{12}R_{23}+I_\lambda R_{13}.
                                                        \tag{G70.9}
\]
There is a localization point to specify in this reuse.
The printed Lemma~8.7 assumes $I_P\subset P R_P$;
(G70.2) alone gives the prime $\mathfrak q$ containing
$I_\lambda$, not that stronger inclusion.
Run its idempotent and Nakayama argument over
$(R_\lambda)_{\mathfrak q}$ modulo $\mathfrak q$ instead.
Its residue field is $\operatorname{Frac}(\Lambda'/P)$,
since $R_\lambda/I_\lambda=\Lambda'/\operatorname{Eis}_\lambda$.
The inertia distinctions proved above give precisely the
required idempotent, and the projection of the first
ordinary filtration piece onto the $\delta_3$-line is
surjective in this local ring.

Put $C=I_1/J_1$.  This module is annihilated by
$I_\lambda$.  It follows that $C_P=C_{\mathfrak q}$:
every element of $R_\lambda/I_\lambda$ is a scalar,
and its inverse is needed exactly when that scalar is
outside $P$.  More explicitly, if a local lift has
denominator $r\notin\mathfrak q$, write
$r\equiv a\pmod{I_\lambda}$ with $a\in\Lambda'\setminus P$.
On $C$, multiplication by $r$ is multiplication by $a$.
Thus the finitely many lift denominators in the ordinary
cocycle calculation can be cleared on this quotient by
a single $t_P\in\Lambda'\setminus P$.
The terms in $R_{12}R_{23}$ vanish on $C$, exactly as in
the proof of Corollary~8.9.  Its global injectivity proof
uses only the residual distinction, and its local
calculation now shows, after this harmless $P$-unit
multiplication, that the same first cocycle has its image
in the indicated Selmer group.  This proves the needed
$P$-localized injection without presuming the printed
stronger inclusion for all of $R_P$.
Theorem~8.11 and $\operatorname{ord}_P L_D^S=0$ then give
$C_P=0$, or equivalently $(I_1)_P=(J_1)_P$.
If $R_{23}(\mathfrak a)=0$, then also
$R_{12}(\mathfrak a)=0$.  After passage to the component
and localization at $\mathfrak q/\mathfrak a$, (G70.9)
says $M=I_\lambda M$ for the finite module coming from
$R_{13}$.  The ideal is inside the actual local maximal
ideal, so Nakayama gives $M=0$.  A nonzero fractional
ideal in a field cannot be killed by localization; hence
$R_{13}(\mathfrak a)=0$.

The trace pairing on this generic component is
nondegenerate: otherwise clear a regular denominator
and obtain a nonzero element of the trace radical that
was already divided out.  A nonzero $(i,j)$ matrix entry
therefore forces a nonzero $(j,i)$ entry by pairing with
$E_{ij}$.  All opposite entries consequently vanish too,
and the generic trace is a sum of three characters.
At an arbitrarily regular classical point on this
component, the precise non-CAP and discrete-packet
argument used on printed p.~101 of Hsieh's proof makes
these characters individually conjugate self-dual.
Reduction contradicts \textup{(Dist)}, because conjugate
duality interchanges the distinct residual endpoint
characters in (G70.3).  These non-CAP and packet results
are retained published inputs, not new elementary claims.
Thus $R_{23}(\mathfrak a)\ne0$.

Symmetry now gives $E_{12},E_{23}$ in the generic image;
trace nondegeneracy gives $E_{21},E_{32}$; products give
$E_{13},E_{31}$.  Together with the diagonal idempotents
these span $M_3$.  Choose nine group elements spanning
that algebra.  Their trace-Gram determinant is nonzero
on the component.  Its nonvanishing open meets the dense
classical points proved above.  At any such point the
nine matrices still span $M_3$, so the attached
representation is absolutely irreducible.  The classical
space being used throughout is the exact
$\lambda$-space, proving the assertion at $q$.
\end{proof}

\paragraph{Precise remaining seed and support gates.}
The second prime $q'$ was introduced only to obtain a
noncyclotomic Euler divisor; it has no prescribed integral
character type in this theorem.  Its new-level components
may or may not have monodromy, and no claim is made in
either direction.  The local automorphic parameter at
$q$ has three distinct tame characters and zero monodromy
by AC-G67-type.  Its identification with the attached
Galois parameter still needs the separately stated
compatibility theorem.  At $269$, when the two endpoint
inertia characters coincide, a nonzero $(1,3)$ monodromy
entry is compatible with those characters, as explicitly
shown in G67.  Thus zero monodromy at every away place
has not been proved.  AC-G64-weight still prevents the
literal classical Hodge-order transport to the symmetric
seed, and AC-G61's actual characteristic-$p$ curve
containment in the Hecke support is not supplied.
There is also no estimate for $q'$ or the enlarged level.

\paragraph{Geometric dictionary and source scope.}
The finite character space here is a tame-cover eigenspace;
the integral trace-one boundary lift and the rational
projector are distinct operations.  The finite-cover
language uses only the already audited SGA~1,
Expos\'e~V, \S\S4 and~7 framework.
The new conclusion is a simple generic local system
on a specified Hecke component, with its exact retained
ordinary and tame types.  No motive, actual symmetric
seed, or unrestricted deformation-ring support theorem
is inferred.  The exact primary inputs are Hsieh
Theorem~4.21, Proposition~7.17, Definition~7.19,
\S8.3 (8.2)--(8.9), \S8.4 Corollary~8.9,
Lemma~8.10, Theorem~8.11 and the non-CAP paragraph in
the proof of Theorem~8.14(ii), together with the
faithful-trace statement of Bella\"iche--Chenevier
Theorem~1.4.4(ii); the original G67 construction retains
its own source hypotheses and audited normalization.
